% TOP_PROBLEM: {"id": 600004, "title": "Baker's Dozen — Periodic hyperbolic outer billiards", "category_id": 6, "category": {"id": 6, "name": "geometry", "display_name": "Geometry", "description": "Euclidean and non-Euclidean geometry, geometric structures.", "slug": "geometry", "order_index": 6, "created_at": "2026-07-31T15:26:25.671Z"}, "status": "open", "problem_number": "AMR-005-0004", "set_id": 13}
% FIRST_POSTED: 2026-10-01
% ENTRY_KIND: statement_only
% UnsolvedMath ID 600004; source catalogue code AMR-005-0004
% !TeX program = lualatex
% Definition and sources only; this file is not an LLM solution attempt.
\documentclass[11pt,a4paper]{article}
\usepackage[margin=25mm]{geometry}
\usepackage{fontspec}
\setmainfont{lmroman10-regular.otf}[BoldFont=lmroman10-bold.otf,
 ItalicFont=lmroman10-italic.otf,BoldItalicFont=lmroman10-bolditalic.otf]
\usepackage{amsmath,amssymb,mathtools}
\usepackage{unicode-math}
\setmathfont{latinmodern-math.otf}
\AtBeginDocument{\let\setminus\smallsetminus}
\usepackage{xurl,hyperref}
\hypersetup{colorlinks=true,urlcolor=blue}
\setcounter{secnumdepth}{0}
\setlength{\parindent}{0pt}
\setlength{\parskip}{5pt}
\setlength{\emergencystretch}{3em}
\begin{document}
\section*{Baker's Dozen — Periodic hyperbolic outer billiards}\label{600004}
{\small UnsolvedMath ID 600004 · AMR-005-0004 · Geometry · AMR Open Problem Lists}
\subsection{Definitions and mathematical statement}
Let $P$ be a compact convex geodesic polygon with nonempty interior in the
oriented hyperbolic plane $\mathbb H^2$. For an exterior point $x$, choose its
right supporting geodesic to $P$. Except when the support contains a side,
it touches a unique vertex $v$. Define $T_P(x)$ by the hyperbolic half-turn
about $v$: the points $x,v,T_P(x)$ lie on a geodesic and $v$ is their midpoint.
The other orientation gives the inverse convention.

Each half-turn extends to the ideal boundary $\partial_\infty\mathbb H^2$.
On that boundary, adjacent sector formulas agree at their common endpoint:
both half-turns interchange the ideal endpoints of the intervening side.
They therefore define a continuous, piecewise projective circle map
$\widehat T_P$ on the entire ideal boundary.

The conjecture is that, for every such $P$, there are an integer $m\ge1$
and either an ordinary exterior point $x$ with $T_P^m(x)=x$, whose orbit
avoids the singular supporting-side rays, or an ideal point $\xi$ with
$\widehat T_P^m(\xi)=\xi$. Every ideal point is permitted, including sector
endpoints. Thus a periodic orbit at infinity is permitted.
No uniform bound on $m$ is requested, and a limiting sequence of almost
periodic trajectories alone does not supply the required orbit.
\subsection{Short English statement}
Must every convex polygon in the hyperbolic plane have a periodic outer
billiard trajectory, allowing the trajectory to lie on the ideal boundary?
\subsection{Sources}
\begin{itemize}
\item[S1] ULAM.AI and the UnsolvedMath Contributors, supplied problems.json, record 600004 (AMR-005-0004), AMR Open Problem Lists. Catalogue identity and source transcription; CC BY 4.0.
\url{https://huggingface.co/datasets/ulamai/UnsolvedMath}.
\item[S2] Tabachnikov - A Baker's Dozen of Problems (2015), Conjecture 2. Original source indexed by AMR-005-0004.
\url{https://amj.math.stonybrook.edu/html-articles/Files-2015-2024/14-01/}.
\end{itemize}
\end{document}
